%% Copyright Roberto Di Cosmo, Jun Furuse, Didier Remy, and Pierre Weis
%% All rights reserved.

%% $Id$

\documentclass[12pt]{article}

\usepackage{graphicx}
\usepackage {pstcol,color,pst-node}
\usepackage{advi}
\usepackage{advi-annot}
\usepackage{advi-graphicx}
\usepackage{demo}
\usepackage{superpose}

\begin{document}

\newpage

\subsection* {{\ActiveDVI} has two components}

\begin{itemize}
\item {\bf a DVI} previewer based on \verb"mldvi" (by Alexandre Miquel)
\begin {itemize}
\item[+] dvips
  \textcolor{c1}{c}%
  \textcolor{c2}{o}%
  \textcolor{c3}{l}%
  \textcolor{c4}{o}%
  \textcolor{c5}{u}%
  \textcolor{c6}{r}
  extension (by Alexandre Miquel)
\item[+] psfile specials with cache (by Jun Furuse)
\item[+] GPIC specials (by Xavier Leroy)
\item[+] Postscript specials (by Didier R{\'{e}}my)
\item[+] hyperrefs (by Didier R{\'{e}}my)
\item[+] active anchors (by Didier R{\'{e}}my)
\item[+] background images (by Roberto Di Cosmo)
\item[+] and a few other gadgets...
\end {itemize}
\end{itemize}

\begin{itemize} \item A set of macros to animate your presentation (by
Jun~Furuse, Didier~R{\'{e}}my, Pierre~Weis, and Roberto Di~Cosmo): the
{\LaTeX} package style \verb"advi.sty".

\item An easy way to launch arbitrary demonstration and/or animation
programs from inside your slides \hbox {(by Jun~Furuse and
Pierre~Weis)}.

\item A set of built-in functions to animate transitions from slides
to slides \hbox {(by Jun~Furuse, Didier R{\'{e}}my, and Pierre Weis)}.

\item Some facilities to resize and move texts within the slides
(by Jun~Furuse, Roberto Di~Cosmo, and Pierre~Weis).

\item Hyperref anchors in the presentation to switch by clicking to
another slide of the show, or even to view other files (either DVI,
text or HTML files).

\item A package to annotate texts (using active anchors).

\end{itemize}

\newpage

%Didier's flashy macro
\def \flash #1{\let \do\leflash \do #1\relax}
\def \leflash #1{\ifx #1\relax \def \do{}\else \def \do {#1\adviwait[0.05]\leflash}\fi \do}

\def\goon{\hbox{{\Large $\Rightarrow$}~Type~space~to~go~on}}

\subsection* {Images in presentations}

{\ActiveDVI} supports \verb+\includegraphics+ with native image formats.

\subsection* {EPS (and other!) graphics inclusion using Camlimages}
\begin{quote}
\verb+\includegraphics[width=0.1\textwidth]{caml.eps}+\\
\verb+\includegraphics[width=0.2\textwidth]{caml.eps}+\\
\verb+\includegraphics[width=0.1\textwidth]{world}+\\

This last command looks for a file named {\tt world} with extensions
out of the ones provided by the directive
\verb+\DeclareGraphicsExtensions+ in the \verb|advi-graphicx| package.
In the present case, it finds \verb+world.jpg+, for which a
bounding box specification must be available in the file
\verb+world.bb+. For this image \verb+world.bb+ is:
\begin{verbatim}
%%BoundingBox: 15 208 580 633
\end{verbatim}

\medskip

\includegraphics[width=0.1\textwidth]{../tex/caml.eps}
\includegraphics[width=0.2\textwidth,height=0.09\textwidth]{../tex/caml.eps}
\framebox{\includegraphics[width=0.1\textwidth]{world}}
\end{quote}

\newpage

\subsection* {Other eye candies for presentation}

\medskip\noindent
\docdef \adviwait\docopt{seconds}
\begin{quote}
Wait for \docid{seconds} or wait for ever if no argument is
provided. Wait means do nothing until the user requests to
continue (or to move to another page).
\end{quote}

\medskip\noindent
\docdef \advirecord \doctt{[play]}\docarg{tag}\docarg{latex code}
\begin{quote}
Computes \docid{latex code} and records its DVI output bound to the tag
\docid{tag}. The DVI output is not displayed during recording, unless the
option \doctt{play} is set.

Embedded recordings are prohibited. 
\end{quote}

\medskip\noindent
\docdef \adviplay\docopt{color}\docarg{tag}
\begin{quote}
Replay the DVI previously recorded as \docid{tag}.
The optional argument changes the color to \docid{color} during replay. 
\end{quote}

\newpage

\subsection* {Eye candies (examples)}

\begin{itemize}
   \item mark the text with a tag name:\\
     $\tt \backslash advirecord\{this\}
      \{\advirecord[play]{this}{Warning~to~the~casual~user}\}$

   \item then change the current text color and replay the tagged text:\\
  $\tt{\backslash}adviplay\{this\}$.
\end{itemize}


\noindent
For instance, to display in red the text associated to the tag
{\tt this} above, you can either pass a color option to the play command:
\verb+\adviplay[red]{this}}+\\
\advirecord[play]{that}{\goon}\adviwait%
\textcolor {c1}{\adviplay {this}\adviplay {that}}\adviwait
{\adviplay{this}\adviplay {that}}\adviwait

\noindent
or place the command in the appropriate color environment:
to display the same text in green, type\\
\verb+\textcolor{green}{\adviplay{this}}+\\
\advirecord[play]{goon}{\goon}\adviwait
\textcolor{green}{\adviplay{this}\adviplay{goon}\adviwait}

\noindent
Using pauses and color settings, you can get
\advirecord[play]{this}{various}
\textcolor{red}{\adviplay{this}\adviplay{goon}}\adviwait
\advirecord[play]{this}{color (and fonts)}
\textcolor{c2}{\adviplay{this}\adviplay{goon}}\adviwait
\advirecord[play]{this}{effects}
\textcolor{c3}{\adviplay{this}\adviplay{goon}}\adviwait
\advirecord[play]{this}{in}
\textcolor{c4}{\adviplay{this}\adviplay{goon}}\adviwait
\advirecord[play]{this}{your}
\textcolor{c5}{\adviplay{this}\adviplay{goon}}\adviwait
\advirecord[play]{this}{texts}
\textcolor{c6}{\adviplay{this}\adviplay{goon}}\adviwait
\advirecord[play]{this}{!}\\
\textcolor{red}{\adviplay{this}\adviplay{goon}}\adviwait
Including some strange flashing effect!
\adviplay{goon}
\textcolor{c6}{\adviplay{this}}\adviwait[\smallpause]%
\textcolor{c1}{\adviplay{this}}\adviwait[\smallpause]%
\textcolor{c2}{\adviplay{this}}\adviwait[\smallpause]%
\textcolor{c3}{\adviplay{this}}\adviwait[\smallpause]%
\textcolor{c4}{\adviplay{this}}\adviwait[\smallpause]%
\textcolor{c5}{\adviplay{this}}\adviwait[\smallpause]%
\textcolor{c6}{\adviplay{this}}\adviwait[\smallpause]%
\textcolor{c1}{\adviplay{this}}\adviwait[\smallpause]%
\textcolor{c2}{\adviplay{this}}\adviwait[\smallpause]%
\textcolor{c3}{\adviplay{this}}\adviwait[\smallpause]%
\textcolor{c4}{\adviplay{this}}\adviwait[\smallpause]%
\textcolor{c5}{\adviplay{this}}\adviwait[\smallpause]%
\textcolor{c6}{\adviplay{this}}\adviwait[\smallpause]%

\newpage

\subsection* {Launching applications}

Any Unix command can be launch during the presentation.
This is convenient to call an external program to draw material in the
{\ActiveDVI} window, or to start a demo from within the presentation.

\bigskip

\noindent
To launch the program {\verb"command"} use in your {\LaTeX} source file\\

\verb+\adviembed{command}+\\

\noindent Additionnally, you may precise the size of the rectangular
window where the application will draw its material. {\ActiveDVI} will
reserve that area for your application.\\

\verb+\adviembed[width=w,height=h]{command}+\\

\noindent
Moreover, {\ActiveDVI} provides two special \TeX macros \verb"!g" and
\verb"!p" that specify the geometry and the (X) window identifier that
the application should use. Hence, if {\verb"command"} has option
{\tt -geometry} to specify its X geometry specification, and option
{\verb"-use"} to specify the X window it should use, then\\

\verb+\adviembed[width=w,height=h]{command -geometry !g -use !p}+\\

\noindent will launch \verb"command" with a geometry computed from the
width and height specification \verb"w" and \verb"h"; \verb"command"
will be display in the window \verb"!p" (normally a sub-window of
{\ActiveDVI}'s window).

\newpage

\subsection* {Embedding applications}

There are two kinds of applications you can launch during the
presentation:
\begin{itemize}
\item {\em embeddable applications} run inside a
subwindow of the main {\ActiveDVI}'s window (hence their window is
stuck to the precise place computed by {\ActiveDVI}),
\item {\em non-embeddable applications} run in a window of
their own (hence their window is unrelated to the {\ActiveDVI}'s window and
may be moved independantly).
\end{itemize}

\subsubsection*{Embeddable applications}

For each application to launch, {\ActiveDVI} allocates a new sub
window of the specified size.  If provided, the identifier \verb"!p"
is replaced by the (X) identifier of that sub window, in order for the
application to use that sub window as its root window. Options to
specify the root window may vary depending of the application at hand
(however, \verb"-use" or \verb"-window" are common to specify this
option). Unfortunately, not every X-application is capable of using a
specified window as its (root) window.  If an application can be
rerooted, it is truely embeddable into {\ActiveDVI} (meaning it is
really a part of the {\ActiveDVI} window). Otherwise the application
is not embeddable.


\subsubsection*{Non embeddable applications}

If the application cannot be rerooted (or if the \verb"!p" parameter
is not specified), {\ActiveDVI} cannot control the placement of the
application: the window manager will handle it. However, {\ActiveDVI}
helps by providing a complete X geometry specification: with width and
height, an extra \verb"+x+y" shift is indicated to precisely access
the place where the embedding should occur in the {\ActiveDVI}
window. Hence, the application appears at the right place even if it
cannot be rerooted and is indeed impossible to embed.

\newpage

\subsubsection*{Launching embeddable applications}

Rerootable applications obeys to the X geometry specifications and can
be properly displayed within the {\ActiveDVI} window (embedding).

Examples:

\begin{itemize}
\item The \verb"display" command (from ImageMagick) can show any image into your
presentation using:\\
\verb|\adviembed[width=3cm,height=3.3cm]|\\
\verb|          {display -geometry !g! -window !p logo.gif}|\\
\adviembed[ephemeral=advilogo,width=3cm, height=3.3260869cm]{display -geometry !g! -window !p advilogo.gif}

\item Similarly, \verb|animate| can display animations with:\\
        \verb|\adviembed[width=2cm,height=2cm]|\\
        \verb|          {animate -geometry !g! -window !p anim.gif}|\\
        \adviembed[ephemeral=animate,width=2cm,height=2cm]{animate -geometry !g! -window !p mmm.anim.gif}
\end{itemize}

\newpage

\subsection*{Launching embeddable applications (cont.)}

\begin{itemize}
  \item Any application which supports embedding options can be launched, as\\
   \verb+\adviembed[width=w,height=h]{prog -geometry !g -use !p}+\\
   For instance, using wish:\\
   \verb|\adviembed[persistent,width=5cm,height=1cm]|\\
   \verb|          {wish ./watch -geometry !g -use !p}|\\
   \adviembed[persistent=wish,width=5cm,height=0.71cm]{wish ./watch -geometry !g -use !p}
\adviwait
  \item Possibly interactive ! (We use CamlTk here)\\
   \verb|\adviembed[persistent=taquin,width=6.5cm,height=9.5cm]|\\
   \verb|          {./taquin -geometry !g -use !p dojoji.gif}|\\
  \adviembed[persistent=taquin,width=8.5cm,height=10cm]{./taquin.sh -geometry !g -use !p dojoji.gif}

\end{itemize}

\newpage

\subsection*{Launching non-embeddable applications}

If an X application does not support the specification of its parent
window, it will be automatically handled by the window manager. Then,
depending on the WM at hand, the placements may or may not obey your
geometry specifications. (In some cases, disabling the ``remember window
positions'' option of your WM may help.)

Applications launched that way will not be properly embedded: they
will appear {\em above} not {\em inside} the {\ActiveDVI}
window. (Hence if you move the {\ActiveDVI} window the applications
will not follow (they may even desappear {\em below} the
{\ActiveDVI}'s window)).

\begin{itemize}
 \item Xeyes using the macro command \verb|embed|:\\
   \verb|\adviembed[width=2cm,height=2cm]|\\
   \verb|    {xeyes -bw 0 -shape -geometry !g}|\\
   \adviembed[width=2cm,height=2cm]{xeyes -bw 0 -shape -geometry !g}

 \item An Xterm running an interesting application:\\
   \verb|\adviembed[width=10cm,height=4cm]|\\
   \verb|    {xterm -geometry 50x7+!x+!y -e ocaml}|\\
% Ne marche pas car la g�ometrie de xterm est en caract�res non en points
%   \adviembed[width=10cm,height=4cm]{xterm -geometry !g -e ocaml}
   \adviembed[width=0.1\textwidth,height=0.02\textheight]{xterm -geometry !g -e ocaml}
\end{itemize}

\newpage

\subsection* {Launching non-embeddable applications (cont.)}

You can add music to your presentation, if you have a music
player available: just embed the music player at the precise place
where you want it to play your sounds.

  \verb|\adviembed[width=1mm,height=1mm]|\\
  \verb|    {mpg123 -q music.mp3}|

\advirecord{play}
 {\adviembed[name=mpg123,width=1mm,height=1mm]{mpg123 -q music.mp3}}
\advirecord{stop}
 {\advikillembed{mpg123}}

\advianchor[click]{play}{Play}
\advianchor[click]{stop}{Stop}

\medskip
\noindent
Here, we added active buttons to start and stop the music (don't forget
to add that kind of gadgets in your shows: it can be very frustrating to have
music automatically launched when you are still speaking! :)

\bigskip

\noindent
In conclusion, the \verb|embed| macro can launch any executable
command during your presentations.

\newpage
 
\begin{center}
{\large\bf I have four words for you!}\\[3mm]

%\adviembed[width=10.62cm,height=7.92cm]{xanim +W!p +Sx!w +Sy!h +Zpe -Zr dancemonkeyboy.avi}
%\adviembed[width=10.62cm,height=7.92cm]{plaympeg -s x!wx!h dancemonkeyboy.mpg}
\adviembed[width=15.93cm,height=12.88cm]{mtvp -z -G=!wx!h -W!p dancemonkeyboy.mpg}
%\adviembed[width=15.93cm,height=12.88cm]{mplayer -geometry !wx!h dancemonkeyboy.mpg}
\end{center}

\newpage

\subsection* {Limitations of Postscript specials}

{\ActiveDVI} recognizes and renders Postscript Specials by calling and telling
\verb"ghostview" to display Postscript on the \verb"advi" window. 
Hence, images drawn using PS tricks will be displayed by \verb"advi". 
However, PS tricks make the assumptions that everything is displayed
by Postscript, which is wrong since \verb"advi" displays most of the
material itself.

\subsubsection* {Most PS-tricks do work}
\begin {itemize}

\item[+] Simple drawings work

\item[+] {\tt\string\SpecialCoor} works, {\em i.e.} commands of the form
{\tt \string \rput \{A\}\{bla bla\}} works where {\tt A} is a node

\item[+] Connections between nodes {\tt \string \ncarc}, {\tt \string \ncarc},
also works.
\end {itemize}

\subsubsection* {Unfortunately, some PS-tricks do not work, yet}

\begin{itemize}
\item[-] Labels over arrows \verb"\Aput", \verb"Bput", etc.
(they change the Postscript coordinates \ldots)

\item[-] {\tt pspicture} (idem, but drawings that are not embedded in
pspictures works)

\item[-] {\ActiveDVI} only recognizes binary overlays (display if 0,
hide otherwise). This is sufficient for the prosper documentclass (see
\hyperlink {overlays}{below}).

\end {itemize}



\subsection*{Known bug with PS tricks}

With some versions of \verb"gv" and/or \verb"X11", the PS tricks are not
correctly synchronized. Some drawings may hide text, or conversely.

% Examples: 

% \Rnode {A}{A} \qquad \Rnode {B}{B}
% \ncline{->}{A}{B} \Aput{L}

\newpage

\subsection* {Overlays}
\hypertarget {overlays}{}

Overlays of the \verb"seminar" mode uses PS tricks, but in a manner that
requires characters to be also printed as Postscript. The \verb"advi" package 
redefines \verb"\overlay" so that it draws only the first overlay (and
ignore others).
Hence, \verb"advi" cannot be used to preview slides using \verb"seminar"
overlays. However, a better package for overlays (build on top of seminar)
is \verb"prosper", which only uses the ability to show or hide some
material, which works well with \verb"advi".

Some reimplementation of \verb"\overlay" for \verb"advi" to work with 
seminar overlays might still be possible...  

\newpage

\section* {Printing your slides: the ignore mode}

The package \verb"advi" can be loaded with the \verb"ignore" option:

\smallskip
{\verb"\usepackage[ignore]{advi}"}
\smallskip

\noindent
The effect is to switch off some \verb"advi" gadgets and to render some others 
in Postscript, if possible. This is useful to generate a DVI file that
can be easily turned into Postscript to print a set of slides as an
approximation of the show that {\ActiveDVI} plays.

\newpage

\subsection* {Hyper links to other files}

\begin {itemize}

\item (Recursive) link to \hyperref{demo.dvi}{}{overlays}{\tt demo.dvi}
is produced by

\qquad\verb"\hyperref{demo.dvi}{}{overlays}{\tt demo.dvi}"

A click opens file \verb"demo.dvi" with a new \verb"advi" process and
jumps to the link \verb"overlays" in that file, as if the command

\qquad\verb"advi -html overlays demo.dvi"

had been typed in a shell.

\item The source file of this demo is \hyperref{demo.tex}{}{}{\tt
demo.tex}. A link to it is produced by

\qquad\verb"\hyperref{demo.tex}{}{}{\tt more demo.tex}"

A click on the link opens the file in a new window, as if 

\qquad\verb"more demo.tex"

had been typed in a shell.

\item Currently, the file name should end with either \verb".dvi",
\verb".txt", \verb".tex", \verb".htm", or \verb".html" extensions, but
in the future this should use Mine man Mailcap to find the type of the
file and the proper previewer for that file.

For the time being, the different viewer commands can be set using an
option on the command line (or including the same option in a
\verb".advirc" file).

For instance, you can compare the rendering of \ActiveDVI's manual

\begin{itemize}

\item in PDF, viewing {\tt ../doc/manual.pdf} with the
\hyperref{../doc/manual.pdf}{}{}{viewer for PDF files} (by default
{\tt xpdf}),
\item in PostScript, viewing {\tt ../doc/manual.ps} with the
\hyperref{../doc/manual.ps}{}{}{\tt viewer for PostScript files} (by
default {\tt gv}),
\item in DVI \hyperref{../doc/manual.dvi}{}{}{as rendered by
\ActiveDVI}.

\end{itemize}


\item A warning is emitted on standard error (Unix {\tt stderr}), if
the file extension is not recognized:

\hyperref{bad-talk.ppt}{}{}{\tt bad-talk.ptt}

or the file is non-existent:

\hyperref{nonexistent.dvi}{}{}{\tt nonexistent.dvi}

\end {itemize}

\newpage

\subsection* {Combination of replays and PS tricks}

Replay command can hide or show text or change the color of text. 
However, Postscript colors are hard-wired by {\TeX}, {\em i.e.} they are not
interpreted in the current color context as font colors are.
\noindent
Replay can still be used to simulate the replay of Postscript colors:
it suffices to superpose two identical figures. For instance, 
$$
\psset {linestyle=none,fillstyle=solid}
\rlap
{\advirecord{tagb}{\ovalnode [fillcolor=green]{A}{\textbf {Postscript}}}}%
\advirecord[play]{taga}{\ovalnode [fillcolor=yellow]{B}{\textbf {Postscript}}}
$$
will  be rendered by
\begin{verbatim}
  \psset {linestyle=none,fillstyle=solid}
  \def \PS {\textbf {Postscript}}
  \rlap {\advirecord{tagb}{\ovalnode [fillcolor=green]{A}{\PS}}}%
  \advirecord[play]{taga}{\ovalnode [fillcolor=yellow]{B}{\PS}}}
\end{verbatim}
\begin {enumerate}

\item Type space to change the text color. 
\adviwait\textcolor{red}{\adviplay{taga}}
\begin{verbatim}
    \textcolor{red}{\adviplay{taga}}
\end{verbatim}

\item Type space to change the background.
\adviwait\textcolor{red}{\adviplay{tagb}}
\begin{verbatim}
    \textcolor{red}{\adviplay{tagb}}
\end{verbatim}

\end {enumerate}

\newpage

\subsection*{Active anchors}

Active anchors are much alike HTML anchors, except that instead of
jumping to a reference anchors (usually another page), they 
replay a recorded piece of DVI when the mouse is over the
anchor area. The page is saved and reset to its normal shape when the mouse
leaves the anchor region.

\noindent
Scratch those stamps with the mouse to see how much you win \verb":)"

$$
\ovalnode [linestyle=none,fillstyle=solid,fillcolor=yellow]{A}
  {\advirecord {A}{\textbf {10000}}}
\advianchor{A}{\adviplay[yellow]{A}}
\hspace {4em}
\def \i{\textcolor{green}{.}}
\ovalnode [linestyle=none,fillstyle=solid,fillcolor=green]{B}
  {\advirecord {B}{\i\textbf{Lost}!\i}}
\advianchor{B}{\adviplay[green]{B}}
$$
The left game was typeset as
{\small
\begin{verbatim}
\ovalnode [linestyle=none,fillstyle=solid,fillcolor=yellow]{A}
  {\advirecord {A}{\textbf {10000}}}
\advianchor{A}{\adviplay[yellow]{A}}
\end{verbatim}}%

\medskip
\noindent
The effect of active anchors can also be 
\advianchor{C}{delocalized}.

\SpecialCoor
\advirecord {C}
  {\ncarc[linecolor=red,arcangleA=50,arcangleB=50,linewidth=3pt]
         {<->}{A}{B}\Aput{Fire!}
}


\noindent
Another form of simple annotations is the overriding of words.

\noindent
Another form of%
\begin{superpose}
\advirecord[play]{a1}{\advianchor{a2}{ simple }}\\
\advirecord{a2}{\adviplay[white]{a1}\colorbox{yellow}{ cool }}
\end{superpose}%
annotations is the overriding of%
\begin{superpose}
\advianchor{7}{ words. }\\
\advirecord{7}{\adviplay[white]{6}\colorbox{white}{{ text. }}}\\
\end{superpose}

\noindent
You can also replace larger pieces of text by other%
\begin{superpose}
\advirecord[play]{1}{pieces \advianchor{2}{of} \advianchor{3}{text}}\\
\advirecord{2}{\colorbox{white}{\begin{tabular}{@{}l@{}}
inner tabulars\\ should be enclosed\\ by parentheses\end{tabular}}}\\
\advirecord{3}{\colorbox{white}{The \qquad End.}}\\
\end{superpose}

\newpage

\subsection*{\adviannot{Active Annotations}(2,1.5){Annotations \\are great!}}

\noindent
\adviannot[p=9em]{Active anchors}(1,1.5)
{\centering Active anchors are even greater!}
can also annotate text as \adviannot[C,bg=green]{here}{here=h+e+r+e},
{\adviannot[c]{for instance}(-1.5,1){\textcolor{red}{for example ?}}}.

\medskip

\noindent
{\let \do\adviemphasize
\do{These} \do{annotations} \do{have} \do{been} \do{typeset} 
\do{respectively} \do{as}\let\do\relax:}

\begin{small}
\begin{verbatim}
\adviannot{Active Annotations}(2,1.5){Annotations \\are great!}

\adviannot[p=9em]{Active anchors}(1,1.5)
   {\centering Active anchors are even greater!}

\adviannot[C,bg=green]{here}{here=h+e+r+e}

\adviannot[c]{for instance}(-1.5,1){\textcolor{red}{for example ?}}
\end{verbatim}
\end{small}

Background of annotations defaults to yellow, centering annotation is
reminiscent to the {\tt array} command: {\tt c} means centered,
{\tt r} right ajusted, and {\tt l} left adjusted. Uppercase variants
{\tt C}, {\tt L}, {\tt R}, have the same meaning, plus implicit math
mode environment.

\subsection*{Note}

The {\setanchor{invisible}\adviannot[bg=red]{commands}(-1.5,1){you found it!}}
\begin{quote}
\hsize \linewidth\setbox0\vbox\bgroup

\verb"\usepackage{advi}"

\verb"\usepackage{advi-annot}"

\egroup
\adviemphasize[lightgreen]{\box0}
\end{quote}
must appear after other packages that uses PS tricks, since it redefines
some macros. The package \verb"advi-annot" is necessary to get
the annotation macros.

\vfil\noindent
\adviemphasize[lightblue]
{\em
Some invisible anchors have been hidden in this page, could you find them?}

\newpage

\subsection*{Superpose Package}

This very small package provides superimposition of boxes, making the
surrounding box large enough to fit all the included boxes. This is convenient
when displaying each included box at different pauses.

\noindent
A kind of {\ActiveDVI} movie
\advirecord{0}{\colorbox{white}
{\begin{superpose}
\advirecord{1}{\textcolor{c2}{can be played}}\\
\advirecord{15}{\textcolor{c1}{changing text at each space}}\\
\advirecord{2}{\textcolor{c3}{you can also use pictures}}\\
\advirecord{31}{\
{\includegraphics[width=0.02\textwidth]{../tex/caml.eps}}}\\
\advirecord{32}{\
{\includegraphics[width=0.02\textwidth]{../tex/caml.eps}}\
{\includegraphics[width=0.03\textwidth]{../tex/caml.eps}}}\\
\advirecord{33}{\
{\includegraphics[width=0.02\textwidth]{../tex/caml.eps}}\
{\includegraphics[width=0.03\textwidth]{../tex/caml.eps}}\
{\includegraphics[width=0.04\textwidth]{../tex/caml.eps}}}\\
\advirecord{34}{\
{\includegraphics[width=0.02\textwidth]{../tex/caml.eps}}\
{\includegraphics[width=0.03\textwidth]{../tex/caml.eps}}\
{\includegraphics[width=0.04\textwidth]{../tex/caml.eps}}\
{\includegraphics[width=0.06\textwidth]{../tex/caml.eps}}}\\
\advirecord{35}{\
{\includegraphics[width=0.02\textwidth]{../tex/caml.eps}}\
{\includegraphics[width=0.03\textwidth]{../tex/caml.eps}}\
{\includegraphics[width=0.04\textwidth]{../tex/caml.eps}}\
{\includegraphics[width=0.06\textwidth]{../tex/caml.eps}}\
{\includegraphics[width=0.08\textwidth]{../tex/caml.eps}}}\\
\advirecord{36}{\
{\includegraphics[width=0.02\textwidth]{../tex/caml.eps}}\
{\includegraphics[width=0.03\textwidth]{../tex/caml.eps}}\
{\includegraphics[width=0.04\textwidth]{../tex/caml.eps}}\
{\includegraphics[width=0.06\textwidth]{../tex/caml.eps}}\
{\includegraphics[width=0.08\textwidth]{../tex/caml.eps}}\
{\includegraphics[width=0.1\textwidth]{../tex/caml.eps}}}\\
\advirecord{37}{{\includegraphics[width=0.1\textwidth]{world}}}\\
\advirecord{4}{{\begin{tabular}{@{}l@{}}
Note that inner tabulars\\ should still be enclosed\\
by parentheses\end{tabular}}}\\
\advirecord{5}{\textcolor{c4}{The End}}
\end{superpose}}}
here.
\adviplay{1}
\adviwait\adviplay{0}\adviplay{15}
\adviwait\adviplay{0}\adviplay{2}
\adviwait\adviplay{0}\adviplay{31}
\adviwait\adviplay{0}\adviplay{32}
\adviwait\adviplay{0}\adviplay{33}
\adviwait\adviplay{0}\adviplay{34}
\adviwait\adviplay{0}\adviplay{35}
\adviwait\adviplay{0}\adviplay{36}
\adviwait\adviplay{0}\adviplay{37}
\adviwait\adviplay{0}\adviplay{4}
\adviwait\adviplay{0}\adviplay{5}

\newpage

\subsection*{Pictures for \TeX}

{\ActiveDVI} can display pictures for \TeX, as generated by
the \verb"gpic" command into a list of \TeX~ \verb"\special" commands.

\bigskip

To obtain a figure, you first write commands for the language
{\em pic} into a file; from this pic file, you can generate a \TeX
file using a {\em pic compiler}, e.g. \verb"gpic"; then include this
compiled file into your presentation, and {\ActiveDVI} will nicely display it.

Let's consider the \verb"pic" source file \verb"trigo.pic" given
\hyperref{trigo_source.dvi}{}{}{here}; we obtain
\verb"trigo.tex" using the command 
\begin{verbatim}
gpic -t trigo.pic > trigo.tex
\end{verbatim}

Then we include \verb"trigo.tex" with the \TeX command
\verb|\expandafter\ifx\csname graph\endcsname\relax \csname newbox\endcsname\graph\fi
\expandafter\ifx\csname graphtemp\endcsname\relax \csname newdimen\endcsname\graphtemp\fi
\setbox\graph=\vtop{\vskip 0pt\hbox{%
    \special{pn 8}%
    \special{ar 1050 1350 900 900 0 6.28319}%
    \graphtemp=.5ex\advance\graphtemp by 1.440in
    \rlap{\kern 0.960in\lower\graphtemp\hbox to 0pt{\hss O\hss}}%
    \graphtemp=.5ex\advance\graphtemp by 2.550in
    \rlap{\kern 1.050in\lower\graphtemp\hbox to 0pt{\hss The trigonometric circle\hss}}%
    \graphtemp=\baselineskip\multiply\graphtemp by 1\divide\graphtemp by 2
    \advance\graphtemp by .5ex\advance\graphtemp by 1.350in
    \rlap{\kern 1.950in\lower\graphtemp\hbox to 0pt{~1\hss}}%
    \graphtemp=\baselineskip\multiply\graphtemp by -1\divide\graphtemp by 2
    \advance\graphtemp by .5ex\advance\graphtemp by 0.450in
    \rlap{\kern 1.050in\lower\graphtemp\hbox to 0pt{\hss 1~}}%
    \special{pa 0 1350}%
    \special{pa 2400 1350}%
    \special{fp}%
    \special{sh 1.000}%
    \special{pa 2300 1325}%
    \special{pa 2400 1350}%
    \special{pa 2300 1375}%
    \special{pa 2300 1325}%
    \special{fp}%
    \graphtemp=\baselineskip\multiply\graphtemp by 1\divide\graphtemp by 2
    \advance\graphtemp by .5ex\advance\graphtemp by 1.350in
    \rlap{\kern 2.400in\lower\graphtemp\hbox to 0pt{\hss $x$}}%
    \special{pa 1050 2400}%
    \special{pa 1050 0}%
    \special{fp}%
    \special{sh 1.000}%
    \special{pa 1025 100}%
    \special{pa 1050 0}%
    \special{pa 1075 100}%
    \special{pa 1025 100}%
    \special{fp}%
    \graphtemp=.5ex\advance\graphtemp by 0.000in
    \rlap{\kern 1.050in\lower\graphtemp\hbox to 0pt{\hss $y$~}}%
    \special{ar 1282 1319 250 250 -0.886486 0.125474}%
    \special{sh 1.000}%
    \special{pa 1502 1208}%
    \special{pa 1440 1125}%
    \special{pa 1533 1169}%
    \special{pa 1502 1208}%
    \special{fp}%
    \graphtemp=\baselineskip\multiply\graphtemp by -1\divide\graphtemp by 2
    \advance\graphtemp by .5ex\advance\graphtemp by 1.319in
    \rlap{\kern 1.532in\lower\graphtemp\hbox to 0pt{~$\theta$\hss}}%
    \special{pa 1050 1350}%
    \special{pa 1829 900}%
    \special{fp}%
    \graphtemp=.5ex\advance\graphtemp by 0.800in
    \rlap{\kern 1.929in\lower\graphtemp\hbox to 0pt{\hss $M$\hss}}%
    \special{pa 1829 900}%
    \special{pa 1050 900}%
    \special{fp}%
    \special{sh 1.000}%
    \special{pa 1150 925}%
    \special{pa 1050 900}%
    \special{pa 1150 875}%
    \special{pa 1150 925}%
    \special{fp}%
    \graphtemp=.5ex\advance\graphtemp by 1.125in
    \rlap{\kern 0.810in\lower\graphtemp\hbox to 0pt{\hss $\sin\left(\theta\right)\left\{\begin{array}{c} \\ \\ \end{array}\right.$\hss}}%
    \special{pa 1829 900}%
    \special{pa 1829 1350}%
    \special{fp}%
    \special{sh 1.000}%
    \special{pa 1854 1250}%
    \special{pa 1829 1350}%
    \special{pa 1804 1250}%
    \special{pa 1854 1250}%
    \special{fp}%
    \graphtemp=.5ex\advance\graphtemp by 1.350in
    \rlap{\kern 1.829in\lower\graphtemp\hbox to 0pt{\hss $\underbrace{\hspace*{1.95cm}}_{\displaystyle\cos~\left(\theta\right)~}$}}%
    \hbox{\vrule depth2.550in width0pt height 0pt}%
    \kern 2.400in
  }%
}%
\showgraph
| and get:

\def\showgraph{%
  \par\medskip\centerline{\raise 1em\box\graph}\bigskip\noindent\ignorespaces}

\expandafter\ifx\csname graph\endcsname\relax \csname newbox\endcsname\graph\fi
\expandafter\ifx\csname graphtemp\endcsname\relax \csname newdimen\endcsname\graphtemp\fi
\setbox\graph=\vtop{\vskip 0pt\hbox{%
    \special{pn 8}%
    \special{ar 1050 1350 900 900 0 6.28319}%
    \graphtemp=.5ex\advance\graphtemp by 1.440in
    \rlap{\kern 0.960in\lower\graphtemp\hbox to 0pt{\hss O\hss}}%
    \graphtemp=.5ex\advance\graphtemp by 2.550in
    \rlap{\kern 1.050in\lower\graphtemp\hbox to 0pt{\hss The trigonometric circle\hss}}%
    \graphtemp=\baselineskip\multiply\graphtemp by 1\divide\graphtemp by 2
    \advance\graphtemp by .5ex\advance\graphtemp by 1.350in
    \rlap{\kern 1.950in\lower\graphtemp\hbox to 0pt{~1\hss}}%
    \graphtemp=\baselineskip\multiply\graphtemp by -1\divide\graphtemp by 2
    \advance\graphtemp by .5ex\advance\graphtemp by 0.450in
    \rlap{\kern 1.050in\lower\graphtemp\hbox to 0pt{\hss 1~}}%
    \special{pa 0 1350}%
    \special{pa 2400 1350}%
    \special{fp}%
    \special{sh 1.000}%
    \special{pa 2300 1325}%
    \special{pa 2400 1350}%
    \special{pa 2300 1375}%
    \special{pa 2300 1325}%
    \special{fp}%
    \graphtemp=\baselineskip\multiply\graphtemp by 1\divide\graphtemp by 2
    \advance\graphtemp by .5ex\advance\graphtemp by 1.350in
    \rlap{\kern 2.400in\lower\graphtemp\hbox to 0pt{\hss $x$}}%
    \special{pa 1050 2400}%
    \special{pa 1050 0}%
    \special{fp}%
    \special{sh 1.000}%
    \special{pa 1025 100}%
    \special{pa 1050 0}%
    \special{pa 1075 100}%
    \special{pa 1025 100}%
    \special{fp}%
    \graphtemp=.5ex\advance\graphtemp by 0.000in
    \rlap{\kern 1.050in\lower\graphtemp\hbox to 0pt{\hss $y$~}}%
    \special{ar 1282 1319 250 250 -0.886486 0.125474}%
    \special{sh 1.000}%
    \special{pa 1502 1208}%
    \special{pa 1440 1125}%
    \special{pa 1533 1169}%
    \special{pa 1502 1208}%
    \special{fp}%
    \graphtemp=\baselineskip\multiply\graphtemp by -1\divide\graphtemp by 2
    \advance\graphtemp by .5ex\advance\graphtemp by 1.319in
    \rlap{\kern 1.532in\lower\graphtemp\hbox to 0pt{~$\theta$\hss}}%
    \special{pa 1050 1350}%
    \special{pa 1829 900}%
    \special{fp}%
    \graphtemp=.5ex\advance\graphtemp by 0.800in
    \rlap{\kern 1.929in\lower\graphtemp\hbox to 0pt{\hss $M$\hss}}%
    \special{pa 1829 900}%
    \special{pa 1050 900}%
    \special{fp}%
    \special{sh 1.000}%
    \special{pa 1150 925}%
    \special{pa 1050 900}%
    \special{pa 1150 875}%
    \special{pa 1150 925}%
    \special{fp}%
    \graphtemp=.5ex\advance\graphtemp by 1.125in
    \rlap{\kern 0.810in\lower\graphtemp\hbox to 0pt{\hss $\sin\left(\theta\right)\left\{\begin{array}{c} \\ \\ \end{array}\right.$\hss}}%
    \special{pa 1829 900}%
    \special{pa 1829 1350}%
    \special{fp}%
    \special{sh 1.000}%
    \special{pa 1854 1250}%
    \special{pa 1829 1350}%
    \special{pa 1804 1250}%
    \special{pa 1854 1250}%
    \special{fp}%
    \graphtemp=.5ex\advance\graphtemp by 1.350in
    \rlap{\kern 1.829in\lower\graphtemp\hbox to 0pt{\hss $\underbrace{\hspace*{1.95cm}}_{\displaystyle\cos~\left(\theta\right)~}$}}%
    \hbox{\vrule depth2.550in width0pt height 0pt}%
    \kern 2.400in
  }%
}%
\showgraph


See the \verb"gpic" command manual for more information about
\verb"gpic" and the programming language \verb"pic".

\lastpage

~\vfill
\begin{center}
\includegraphics[width=0.5\textwidth]{../tex/advilogo.eps}\\
{\Large \bf That's all, folks!!}
\end{center}
\vfill

\end{document}
